\chapter{第十八讲\quad 人身犯罪}
\label{lecture:18}
\noindent{\small\hyperref[ref:lecture:18]{查阅本讲扩展规则}}\par

本讲把同一外观行为按时间拆成「行为、对象、目的、结果」四列。$a$ 为被评价的行为，$r$ 为发生的结果，$C(a,r)$ 表示结果能够归属于该行为。下列蕴涵均以责任年龄、责任能力等共同条件已经满足为前提，涉及未证实的条件仍须另行论证。

\section{第一节\quad 侵犯生命、身体的犯罪}

\paragraph{法条与模型}
依《刑法》第232、233、234、235条分别检查死亡故意、伤害故意、死亡过失和重伤过失。令 $I_D(a)$ 表示对死亡希望或放任，$I_H(a)$ 表示伤害故意，$N_D(a)$ 表示对死亡存在过失，$D(r)$ 表示死亡。把能够归责的死亡写成 $R_D(a):=\exists r[D(r)\land C(a,r)]$。在其他条件固定时，
\[
\begin{aligned}
 I_D(a)\land R_D(a)&\Longrightarrow \text{故意杀人既遂},\\
 I_H(a)\land\neg I_D(a)\land N_D(a)\land R_D(a)
   &\Longrightarrow \text{故意伤害致死},\\
 N_D(a)\land\neg I_H(a)\land\neg I_D(a)\land R_D(a)
   &\Longrightarrow \text{审查过失致人死亡}.
\end{aligned}
\]
最后一行还要排查交通肇事等特别规定。伤害故意根据下手部位、力度、工具、言行等证据认定。由 $R_D(a)$ 单独推不出 $I_D(a)$。直接故意杀人已经着手，因外力未得逞，另按第23条审查未遂。

\paragraph{推导}
甲为使乙受伤实施殴打，题设明确甲排斥死亡结果，但应当预见剧烈殴打可能致死。乙因该殴打死亡，且无阻却事由。于是 $I_H,\neg I_D,N_D,R_D$ 四项成立，由合取消去其他分支，再由第二行得到故意伤害致死。若换成乙躲避后死于与殴打无关的独立事故，则 $R_D$ 缺失，死亡加重项随之失去，仍应按原行为及其实际伤害评价。

\paragraph{器官与扶养}
第234条之一把组织出售人体器官独立规定。区分 $O$（组织出卖）、$U$（活体无同意、未满十八周岁，或强迫、欺骗捐献）和 $S$（符合本条第三款的尸体器官情形）。仅有购买事实推不出 $O$。$U$ 触发第234条或第232条的特别指向，死亡及罪过仍分开分析。$S$ 则指向第302条。题设成年人同意捐献给特定亲属，组织者用虚假事实使其为别人捐献，已证实欺骗捐献，故先进入 $U$ 的伤害或杀人评价，不可把外观同意直接记为有效承诺。

第261条的遗弃路径为「无独立生活能力 $\land$ 扶养义务 $\land$ 有能力而拒绝扶养 $\land$ 情节恶劣」。改按不作为杀人须另外证明保证人义务、救助可能、结果归责、与作为杀人的等价性及杀人故意。一个有扶养义务且长期拒绝照护的案件，即使情节恶劣，也可能没有已经迫近且可由其救助避免的死亡危险。遗弃成立不能推出不作为杀人。

\section{第二节\quad 侵犯性权利的犯罪}

\paragraph{法条与分流}
设 $u$ 为被害人案发时周岁，$F$ 为女性对象，$K$ 为行为人对幼女身份的法定明知，$X$ 为第236条所评价的性行为，$W$ 为暴力、胁迫等手段侵害性自主。$X_c$ 为与该未成年女性发生第236条之一所评价的性关系，$H$ 为对该对象负有特殊职责。$M$ 为猥亵行为，$K_c$ 为相应犯罪所需的儿童身份认识。用四个入口整理第236条、第236条之一和第237条：
\[
\begin{aligned}
 F\land X\land W&\Longrightarrow \text{审查强奸罪的一般入口},\\
 F\land X\land (u<14)\land K&\Longrightarrow \text{幼女特别入口},\\
 F\land X_c\land (14\le u<16)\land H&\Longrightarrow \text{照护职责入口},\\
 M\land (u<14)\land K_c&\Longrightarrow \text{审查猥亵儿童罪}.
\end{aligned}
\]
第237条第一款的强制猥亵他人或者侮辱妇女，另检查强制手段和具体对象。猥亵儿童有独立条款。2023年四部门办案意见第17条对 $K$ 分段认定。对不满十二周岁的幼女实施性侵害，依法认定明知。对已满十二、不满十四的，结合发育状况、言谈举止、衣着特征、生活作息等，核对行为人是否能够观察出对方可能是幼女。$K_c$ 仍按相应对象和证据独立审查。一般强奸路径和照护职责路径同时满足时，按第236条之一第二款择处罚较重的规定。

\paragraph{推导}
甲对十五周岁的女性乙负有明确监护职责，与乙发生第236条之一所指性关系，题设无证据支持强制手段。由 $14\le 15<16$ 得年龄条件，由监护得 $H$，合取 $F,X_c$ 得照护职责入口。缺少 $W$ 的证据不能写作已经证明 $\neg W$，只能说现有事实未推出一般强奸入口。若改为甲利用优势地位迫使乙发生第236条所涉行为，则补入 $X,W$，两条路径均需比对。

\paragraph{边界}
年龄、在场多人、共同正犯关系、重伤死亡是不同坐标。比如未成年被害案件中的「公共场所当众」，2023年办案意见第18条要求公共场所和其他多人在场，未要求每个人都看见。另起的灭口杀害若与前罪的实行行为无归责联系，不能直接塞入强奸致死。对不满十四周岁从重和不满十周岁等情形升格，也要分别依据第236条相应款项，避免把从重写成升档。

\section{第三节\quad 侵犯人身自由的犯罪}

\paragraph{法条与教材模型}
定义 $L$ 为非法剥夺人身自由，$B$ 为控制人质时已经具有向第三人提出不法要求的目的，$S$ 为出卖目的，$A$ 为第240条第二款列举行为之一，$Q$ 为拐骗未满十四周岁者脱离家庭或监护。先把事实分别代入第238、239、240、241、262条，再处理重叠。
\[
\begin{aligned}
 L&\Longrightarrow \text{审查非法拘禁},\\
 L\land B&\Longrightarrow \text{审查绑架},\\
 S\land A&\Longrightarrow \text{审查拐卖妇女、儿童},\\
 Q&\Longrightarrow \text{审查拐骗儿童}.
\end{aligned}
\]
收买被拐卖的妇女、儿童按第241条评价。收买时本有出卖目的，或收买后又出卖，分别审查第240条和第241条第五款的指向。收买后另行强奸、拘禁等，依第241条的并罚规则处理。以勒索财物目的偷盗婴幼儿，由第239条第三款专门接入。

\paragraph{推导}
甲控制乙作为人质，着手控制时即计划逼乙的家属付款，之后尚未向家属发出信息即被制止。控制事实给出 $L$，目的及时间给出 $B$。按教材所列多数说，取得实力控制即完成绑架的实行行为，所以即使勒索行为尚未开始，仍能推出绑架既遂。教材所列少数说把对第三人的要求纳入复合实行行为，会改变既未遂结论。争点应写在选定的规则上，案情集合保持一致。若目的只是向乙当场强取乙的钱，则 $B$ 缺失，改查抢劫的强制取财结构。若为索取确有的债务而非法扣押债务人，则先适用第238条第三款的特别路径。

\paragraph{结果与竞合}
非法拘禁致重伤、死亡先核对第238条第二款的结果加重入口，再查有无该款后句的使用暴力致伤残、死亡情形。拘禁者仅造成一般殴打时，不可跳到暴力致伤残、死亡分支。绑架后对人质的杀害、故意伤害致重伤死亡依第239条第二款检验。第240条列举的奸淫被拐卖妇女和造成特定对象重伤死亡等也有专门升格入口，应先检验专条再讨论数罪。

\section{第四节\quad 侵犯名誉权、人身权的犯罪}

\paragraph{法条与谓词}
令 $F$ 为捏造事实，$P$ 为向外传播或公然实施，$G$ 为使他人受刑事追究的目的，$J$ 为司法工作人员身份，$V$ 为暴力手段。第246条审查公然贬损或捏造事实诽谤及情节严重。第243条则审查捏造事实诬告、$G$ 及情节严重。第247条可拆为两个不同对象的命题：
\[
\begin{aligned}
 J\land\text{犯罪嫌疑人或被告人}\land\text{刑讯逼供}
   &\Longrightarrow \text{刑讯逼供入口},\\
 J\land\text{证人}\land V\land\text{逼取证言}
   &\Longrightarrow \text{暴力取证入口}.
\end{aligned}
\]

\paragraph{推导}
甲伪造乙涉嫌犯罪的证据并向侦查机关诬告，题设明确甲意在使乙被刑事追究且情节严重。由伪造得 $F$，由递交机关和明确目的得诬告及 $G$，经第243条进入诬告陷害。若甲仅在群聊中散布诋毁乙的虚假消息，已证明其只有贬损名誉目的，则缺 $G$，转查第246条。若只是错告、检举失实，按第243条第三款排除该条的故意诬陷路径。

\paragraph{边界与补充入口}
监管人员殴打、体罚虐待被监管人且情节严重按第248条。非法搜查、非法侵入住宅按第245条。违反国家有关规定出售、提供或者非法获取个人信息并达到规定情节，按第253条之一。第244条检查强制劳动手段，明知而协助有专门条款。第244条之一检查未满十六周岁、危重劳动和情节严重。故「招工对象年龄十五岁」单独不能推出危重劳动罪，作业环境与情节须另证。第246条一般告诉才处理，其法定例外是程序层判断，不能用于反推诽谤要件。

\section{第五节\quad 侵犯婚姻家庭的犯罪}

\paragraph{法条与模型}
把 $M$（有配偶）、$K$（明知对方有配偶）、$A$（结婚或依适用解释作同类评价的事实）、$S$（明知对方是现役军人的配偶）、$C$（同居）分开。第258条重婚的两入口是 $(M\land A)$ 和 $(K\land A)$。第259条第一款破坏军婚的入口是 $S\land(C\lor A)$。这里的「同居」按刑法语境审查持续稳定共同生活，婚姻民事效力另行评价。第259条第二款另规定利用职权、从属关系，以胁迫手段奸淫现役军人妻子的处理。

\paragraph{推导}
甲无配偶，明知乙是现役军人的配偶，并长期与乙稳定共同生活，题设明确双方没有办理婚姻登记，也不存在以夫妻名义共同生活的证据。事实给出 $S,C$，由第259条进入破坏军婚。对重婚，现有证据仍未建立 $A$，不能以同居径行推出该罪。如果乙只是退役军人的配偶，则 $S$ 缺失，重婚要件和其他法律后果仍应独立审查。

\paragraph{暴力与虐待}
第257条要求暴力干涉婚姻自由，并分离基本犯的告诉规则和致死入口。第260条的虐待以家庭成员关系和情节恶劣等为条件。第260条之一则检查监护、看护职责和特定对象，单位也可能成为主体，并有择重规则。若照护机构人员对被照护老人实施情节恶劣的持续虐待，由照护职责和对象可进入第260条之一。二人并无家庭成员关系时，不能把机构照护替代第260条的身份条件。

\section*{闭卷记忆与验算}
\paragraph{复述路线}
先画五格「生、性、自由、名身、婚家」。每格先写对象与身份，再问行为的即时目的，再查可归责结果和特别法条。压缩成「对象先分，身份复核，目的同刻，结果另证，专条收尾」。其中「同刻」对应 $B,S$ 的发生时点，「另证」对应 $C(a,r)$ 和相应结果罪过。

\begin{enumerate}
\item 题：甲先为催还真实债务非法拘禁债务人，之后才萌生向其家属索取无关财物的想法。能否把最初行为直接记作绑架？答：不能。初段先按第238条第三款。后段另查该目的形成后有无继续以人质手段控制他人等行为，不能把后来的 $B$ 回填到初始时点。
\item 题：甲以伤害故意殴打乙，乙死亡。还缺哪两类事实才能判定故意伤害致死？答：一是死亡能归责于该殴打，二是甲对死亡有过失，并排查是否其实放任死亡。死亡结果本身不提供这两类事实。
\item 题：十五周岁的被害人对照护者表面表示同意，是否可以直接排除犯罪？答：先查年龄、特殊职责和行为三项。对第236条之一，表面同意不能单独排除。再检查有无利用优势地位迫使等一般强奸的事实。
\end{enumerate}
